\section{Independent proof by moving jobs with \texorpdfstring{$p=u$}{p = u} to a prefix}
\label{app:cap-bubbling}
For completeness, we give a second proof of the pairwise estimate behind
\cref{thm:opt:optional-middle-upper}.  This
argument is also a useful audit of the endpoint type $U$, which represents
$p=u$, in the five-endpoint reduction \cref{eq:opt:five-endpoints}.  In
particular, it never substitutes $p=u$ into the formula for type $Q$, whose
value is just above the current threshold and strictly below $u$, before the
$U$ endpoints have actually been moved.

Recall the endpoint symbols from \cref{eq:opt:five-endpoints}: $Z$ has value
zero; $E$ is a positive value tending to zero; $M$ equals the current
threshold and is processed after its test; $Q$ lies just above that threshold
and is placed in the final queue; and $U$ has value $p=u$.  A word in these
symbols records their reveal order.
For neighboring symbols $X,U$, the expression $UX-XU$ denotes the change in
the pair objective when $U$ is moved one position to the left.

Work first with class sizes divided by the total number of jobs; we call
these normalized sizes their fluid masses.  Immediately before neighboring
symbols
$XU$, let $E_0$ be the preceding $E$-mass, $P$ the preceding mass of types
$M,Q,U$, $N_{\mathrm{def}}$ its deferred mass of types $Q,U$, $x$ the
remaining mass, and $w$ the mass of $U$-endpoints strictly after the pair.
Put
\begin{equation}
 y=\frac{N_{\mathrm{def}}-R(P+E_0)}x,
 \qquad
 \eta=\frac{N_{\mathrm{def}}-RP}{x}=y+R\frac{E_0}{x}\le0.
 \label{eq:opt:bubble-state}
\end{equation}
Direct differentiation of the exact pair objective
\eqref{eq:opt:pair-objective}, including the mutual pair and all future
$U$-endpoints, gives the two nontrivial swap differences
\begin{align}
 UM-MU
 &=1+\left(1+\eta-\frac{Rw}{x}\right)
       (y-c)a_c'(y),
 \label{eq:opt:cap-M-commutator}\\
 UQ-QU
 &=\left(\eta-\frac{cw}{x}\right)
       (y-c)a_c'(y).
 \label{eq:opt:cap-Q-commutator}
\end{align}
The second quantity is nonnegative.  For the first, if
$B=1+\eta-Rw/x\le0$, its second term is nonnegative.  If $B>0$, then
$B\le1$, and on $-1<y\le0$
\[
 \frac{(c-y)B}{c-2y}\le1.
\]
Thus \eqref{eq:opt:cap-M-commutator} is nonnegative as well.  Below
$y=-1$, the threshold is flat.  Direct evaluation also gives
\begin{equation}
 UZ-ZU=UE-EU=1.
 \label{eq:opt:cap-zero-commutators}
\end{equation}
The terms involving the later $U$-mass $w$ are crucial: they make every
neighboring swap valid even while other $U$-endpoints remain to the right.
Repeatedly moving $U$ left in this way---the ``bubbling'' in the section
title---puts all $U$-endpoints into one prefix.

After that prefix there is no future $U$-endpoint.  Let $e$ be the preceding
$E$-mass,
$b$ the preceding mass of types $M,Q$, and $d$ the preceding deferred
$Q$-mass, and define
\[
 y=\frac{d-R(e+b)}x,
 \quad v=\frac ex,
 \quad h=d-Rb,
 \quad k=x+d-Rb.
\]
The four swap differences involving only the types $Z,E,M,Q$ are
\begin{align}
 ME-EM&=a_c+a_c'(R-y)\frac{k}{x},
 &QE-EQ&=1-a_c'(R-y)(-y-Rv),
 \notag\\
 MZ-ZM&=a_c-a_c'y\frac{k}{x},
 &QZ-ZQ&=1-a_c'y\frac{h}{x}.
 \label{eq:opt:ordinary-commutators}
\end{align}
They are nonnegative.  Indeed, in the varying-threshold region
$-1\le y\le0$,
$h/x=y+Rv\in[y,0]$ and $k/x\in[0,1]$.  With $t=-y\in[0,1]$,
\begin{equation}
 a_c'(R-y)t=\frac{t(R+t)}{c+2t}\le1,
 \qquad
 a_c't^2=\frac{t^2}{c+2t}\le1;
 \label{eq:opt:ordinary-commutator-signs}
\end{equation}
the first inequality is equivalent to
$(t-1)(t+c)\le0$.  In the constant-threshold region $y\le-1$, one has
$a_c'=0$.  Hence $Z,E$ can be moved to the end, and a worst word has the form
\begin{equation}
 U^\mu\{M,Q\}^*\{Z,E\}^*.
 \label{eq:opt:canonical-word}
\end{equation}
Here $U^\mu$ denotes a block of $U$-mass $\mu$, and the star denotes an
arbitrary finite sequence of symbols from the set inside the braces.

On the segment containing the types $M,Q$, use the simpler potential
\[
 \widehat W(x,N,N_{\mathrm{def}})=xN_{\mathrm{def}}+x^2f_c\!\left(\frac{N_{\mathrm{def}}-RN}{x}\right).
\]
The identity
\begin{equation}
 2f_c-yf_c'
 =\max\{0,-Rf_c',-Rf_c'+a_c(1+y),
              1-cf_c'+a_cy\}
 \label{eq:opt:ordinary-HJB}
\end{equation}
is a direct rewriting of \eqref{eq:opt:f-hjb-identities}; it bounds the
entire continuation containing only the types $Z,E,M,Q$.  If the prefix of
$U$-endpoints has mass $\mu$, then
after it $x=1-\mu$, $N=N_{\mathrm{def}}=\mu$, and
$y=-c\mu/(1-\mu)$.  Its contribution from pairs of $U$-endpoints plus its worst
continuation is therefore
\begin{equation}
 F_{c,u}(\mu)
 =\frac{2-c(u-1)}2\mu^2+\mu(1-\mu)
 +(1-\mu)^2f_c\!\left(-\frac{c\mu}{1-\mu}\right).
 \label{eq:opt:prefix-objective}
\end{equation}
For $0<\mu<1/R$, differentiation gives
\begin{equation}
 F_{c,u}'(\mu)
 =c\mu\,[Z(c)-u-\operatorname{artanh}\mu].
 \label{eq:opt:prefix-derivative}
\end{equation}
For parameters $(c,m,u)$ satisfying
\cref{eq:opt:cm-parametrization}, the unique maximizer on
$0<\mu<1/R$ is therefore $\mu=m$, and
\begin{equation}
 F_{c,u}(m)
 =f_c(0)+\frac c2(\operatorname{artanh}m-m)=\frac c2.
 \label{eq:opt:prefix-maximum}
\end{equation}
Here $m\le1/R$, since $c,m\le1/\varphi$.  On the flat part
$\mu\ge1/R$,
\begin{equation}
 F_{c,u}(\mu)=\mu-\frac{c(u-1)}2\mu^2.
 \label{eq:opt:prefix-flat}
\end{equation}
Its derivative at $1/R$ is
\[
 \frac{c(1+m)-1}{R}\le0,
\]
and only decreases thereafter.  Thus $F_{c,u}(\mu)\le c/2$ on the whole
unit interval.  At $c=r,m=0,u=\uDet{5}$,
\cref{eq:opt:prefix-derivative,eq:opt:prefix-flat} show that the maximum is
at $\mu=0$, and increasing $u$ only lowers
\eqref{eq:opt:prefix-objective}.

To make this independent exchange proof finite, fix
$\eps>0$ and first replace the last
$\eps n+O(1)$ values by actual zeros.  Earlier states and
actions are unchanged, and only $O(\eps n^2+n)$ pair terms are
affected.  All swaps between $U$ and non-$U$ endpoints among the retained
jobs then
occur with $x\ge\eps$; their displayed normalized swap differences have
uniform Taylor remainders there.  More precisely, after assigning mass
$1/n$ to every job, a normalized neighboring swap has remainder
$O_{u,\eps}(n^{-3})$; there are at most $n^2$ exchanges.  At the
boundary $y=-1$, split a swap at the boundary and use the two one-sided
formulae, whose values agree.  The artificial zeros already form the
last block, so no $U$-endpoint remaining later in the word invalidates
\eqref{eq:opt:ordinary-commutators}.  The normalized one-job remainder in
the change of the potential for the final $M,Q,Z,E$ segment is
$O_{u,\eps}(n^{-2})$.  For fixed $\eps$, first let
$n\to\infty$, and then let $\eps\downarrow0$.  Equations
\eqref{eq:opt:ordinary-HJB}--\eqref{eq:opt:prefix-maximum} give
\[
 \Gamma_n\le c\binom n2+o(n^2),
\]
which independently proves the required size-asymptotic inequalities in
\cref{thm:opt:optional-middle-upper}, both for
$\uDet{4}\le u\le\uDet{5}$ and for $u\ge\uDet{5}$.
